\documentclass[11pt,a4paper]{article}
\usepackage{fontspec}
\usepackage{kotex}
\usepackage{geometry}
\usepackage{graphicx}
\usepackage{xcolor}
\usepackage{tcolorbox}
\usepackage{booktabs}
\usepackage{array}
\usepackage{longtable}
\usepackage{colortbl}
\usepackage{fancyhdr}
\usepackage{titlesec}
\usepackage{hyperref}
\usepackage{emoji}
\usepackage{amsmath}
\usepackage{amssymb}
\usepackage{enumitem}

\geometry{left=25mm, right=25mm, top=30mm, bottom=30mm}

\setmainfont{Noto Sans CJK KR}
\setsansfont{Noto Sans CJK KR}
\setmonofont{Noto Sans CJK KR}
\setemojifont{Noto Color Emoji}

\definecolor{primary}{HTML}{1976D2}
\definecolor{darkbrown}{HTML}{3E2723}
\definecolor{mediumbrown}{HTML}{5D4037}
\definecolor{lightbrown}{HTML}{8D6E63}
\definecolor{cream}{HTML}{FFF3E0}
\definecolor{lightgray}{HTML}{F5F5F5}
\definecolor{tableheader}{HTML}{E8E8E8}

\titleformat{\section}
  {\Large\bfseries\color{primary}}
  {\thesection.}{0.5em}{}
\titleformat{\subsection}
  {\large\bfseries\color{darkbrown}}
  {\thesubsection}{0.5em}{}
\titleformat{\subsubsection}
  {\normalsize\bfseries\color{mediumbrown}}
  {\thesubsubsection}{0.5em}{}

\renewcommand{\contentsname}{목차}
\renewcommand{\figurename}{그림}
\renewcommand{\tablename}{표}

\tcbset{
  tipbox/.style={colback=cream,colframe=primary,boxrule=1.5pt,arc=3mm,
    left=5mm,right=5mm,top=3mm,bottom=3mm,fonttitle=\bfseries\color{primary}},
  formulabox/.style={colback=lightgray,colframe=lightgray,boxrule=0pt,arc=2mm,
    left=5mm,right=5mm,top=3mm,bottom=3mm},
  summarybox/.style={colback=white,colframe=primary,boxrule=2pt,arc=4mm,
    left=5mm,right=5mm,top=5mm,bottom=5mm}
}

\hypersetup{colorlinks=true,linkcolor=primary,urlcolor=primary,bookmarks=true,unicode=true}


\pagestyle{fancy}
\fancyhf{}
\fancyhead[R]{\small\color{lightbrown}수업 집중도 분석 프로그램 활용가이드 v1.0.0.0}
\fancyfoot[C]{\thepage}
\renewcommand{\headrulewidth}{0.4pt}

\begin{document}

\begin{titlepage}
\centering
\vspace*{3cm}
{\Huge\bfseries\color{primary}\emoji{bar-chart} 수업 집중도 분석 프로그램}\\[0.5cm]
{\Huge\bfseries\color{primary} 활용가이드}
\\[1.5cm]
{\Large 시뮬레이터 버전: v2.4.0.0}
\\[2cm]
\begin{tcolorbox}[summarybox]
\centering
{\large 이 활용가이드는 수업 집중도 분석 프로그램의 수업 적용 방법을 안내합니다.}
\end{tcolorbox}
\vfill
{\large 사물인터넷과 빅데이터}\\[0.3cm]
{\large 청주교육대학교 컴퓨터교육과}
\vspace{1cm}
\end{titlepage}

\section*{1. 수업 개요}

\begin{center}
\renewcommand{\arraystretch}{1.4}
\begin{tabular}{|p{2.8cm}|p{12.2cm}|}
\hline
\rowcolor{tableheader}
\textbf{항목} & \textbf{내용} \\
\hline
강좌 & 사물인터넷과 빅데이터 \\
\hline
적용 주차 & 9\textasciitilde{}13주차 \\
\hline
자료 역할 & 참고자료 \\
\hline
대상 & 4학년 (프로그래밍 비전공, 초등 예비교사) \\
\hline
수업 시간 & 45분 (1차시) \\
\hline
\end{tabular}
\end{center}

\subsection*{1.1 학습 목표}

$\bullet$ 데이터에서 상관관계를 발견하고 인과관계와 구분할 수 있다

$\bullet$ 규칙 기반 모델과 회귀 모델의 차이점을 비교할 수 있다

$\bullet$ AI가 '판단'이 아닌 '계산'을 수행한다는 점을 설명할 수 있다



\section*{2. 수업 시나리오 (45분)}

\begin{center}
\renewcommand{\arraystretch}{1.4}
\begin{tabular}{|p{1.5cm}|p{1.5cm}|p{7.8cm}|p{4.2cm}|}
\hline
\rowcolor{tableheader}
\textbf{단계} & \textbf{시간} & \textbf{활동} & \textbf{시뮬레이터 활용} \\
\hline
도입 & 5분 & "시끄러운 교실에서 집중이 잘 되나요?" 경험 공유 & — \\
\hline
데이터 & 10분 & 샘플 데이터 로드 후 분석 탭에서 산점도 관찰 & \emoji{inbox-tray} 수집, \emoji{bar-chart} 분석 \\
\hline
실습1 & 10분 & 학생 모델 탭에서 나만의 온도·소음 기준 설정 & \emoji{man-student} 학생 모델 \\
\hline
시연 & 10분 & 교수 모델 학습 시작 \textrightarrow{} 수식·R² 확인 \textrightarrow{} 비교 탭 & \emoji{man-teacher} 교수, \emoji{balance-scale} 비교 \\
\hline
토론 & 10분 & "AI가 집중도를 판단하는 걸까?" 토론 + 정리 & — \\
\hline
\end{tabular}
\end{center}


\section*{3. 학습 질문}

\subsection*{초급 (이해)}

$\bullet$ Q: 온도와 집중도 사이에 어떤 관계가 보이나요?

$\bullet$ Q: 학생 모델과 교수 모델의 가장 큰 차이는 무엇인가요?


\subsection*{중급 (적용)}

$\bullet$ Q: 슬라이더로 온도 기준을 바꾸면 정확도가 어떻게 변하나요?

$\bullet$ Q: R²가 0.3이면 무엇을 의미하나요?


\subsection*{고급 (분석)}

$\bullet$ Q: 상관관계가 있다고 인과관계가 있다고 할 수 있나요? 예를 들어 설명해보세요.

$\bullet$ Q: AI가 '집중도를 판단한다'는 표현이 정확한가요?


\subsection*{확장 (창의)}

$\bullet$ Q: 집중도 외에 이 분석 방법을 적용할 수 있는 다른 분야는 무엇이 있을까요?

$\bullet$ Q: 더 정확한 모델을 만들려면 어떤 데이터가 추가로 필요할까요?



\section*{4. 학생 활동지}

\subsection*{활동: 나만의 집중도 규칙 만들기}

\textbf{목표}: 온도·소음 기준을 설정하여 집중도 예측 모델을 만든다


\begin{center}
\renewcommand{\arraystretch}{1.4}
\begin{tabular}{|p{1.8cm}|p{6.4cm}|p{4.1cm}|p{2.7cm}|}
\hline
\rowcolor{tableheader}
\textbf{시도} & \textbf{최적 온도 범위} & \textbf{소음 기준} & \textbf{정확도} \\
\hline
1 & \_\_\_°C \textasciitilde{} \_\_\_°C & \_\_\_dB & \_\_\_\% \\
\hline
2 & \_\_\_°C \textasciitilde{} \_\_\_°C & \_\_\_dB & \_\_\_\% \\
\hline
3 & \_\_\_°C \textasciitilde{} \_\_\_°C & \_\_\_dB & \_\_\_\% \\
\hline
\end{tabular}
\end{center}

\textbf{비교}: 교수 모델 정확도: \_\_\_\% / 나의 최고 정확도: \_\_\_\%


\textbf{생각해보기}: 두 모델의 차이가 생기는 이유는 무엇인가요?



\section*{5. 평가 관찰 체크리스트}

\begin{center}
\renewcommand{\arraystretch}{1.4}
\begin{tabular}{|p{3.1cm}|p{5.5cm}|p{4.4cm}|p{2.0cm}|}
\hline
\rowcolor{tableheader}
\textbf{평가 항목} & \textbf{상} & \textbf{중} & \textbf{하} \\
\hline
상관 vs 인과 구분 & 예시를 들어 정확히 구분 & 차이를 대략 이해 & 혼동 \\
\hline
모델 비교 & 두 모델 차이를 분석적으로 비교 & 정확도 비교만 수행 & 비교 미수행 \\
\hline
AI 본질 이해 & '계산'과 '판단'의 차이 설명 & AI 작동 원리를 대략 이해 & 이해 부족 \\
\hline
\end{tabular}
\end{center}


\section*{6. 수업 팁}

$\bullet$ \emoji{light-bulb} 샘플 데이터로 시작한 뒤 랜덤 생성으로 확장하면 부담이 줄어듭니다

$\bullet$ \emoji{light-bulb} 학생 모델에서 '틀린 규칙'도 가치가 있음을 강조하세요 — 시행착오가 학습입니다

$\bullet$ \emoji{light-bulb} "AI가 판단한다"는 표현을 학생들이 쓸 때마다 "AI가 계산한다"로 교정하세요



\section*{7. 관련 자료}

\begin{center}
\renewcommand{\arraystretch}{1.4}
\begin{tabular}{|p{1.5cm}|p{6.9cm}|p{6.6cm}|}
\hline
\rowcolor{tableheader}
\textbf{\#} & \textbf{자료명} & \textbf{비고} \\
\hline
1 & 수업집중도분석프로그램\_퀵가이드 & 소프트웨어 사용 중 빠른 참조 \\
\hline
2 & 수업집중도분석프로그램\_사용설명서 & 전체 기능 상세 \\
\hline
3 & 수업집중도분석프로그램\_이론해설서 & 배경 이론 \\
\hline
4 & 수업집중도분석프로그램\_개발참조 & 기술 사양 \\
\hline
\end{tabular}
\end{center}


\vfill
\begin{center}
\small\color{lightbrown}
\textit{수업 집중도 분석 프로그램 활용가이드}
\end{center}
\end{document}